% 行星（综述）
% license CCBYSA3
% type Wiki

本文根据 CC-BY-SA 协议转载翻译自维基百科\href{https://en.wikipedia.org/wiki/Planet}{相关文章}。

\begin{figure}[ht]
\centering
\includegraphics[width=8cm]{./figures/c14ce525d8184ebd.png}
\caption{太阳系八大行星的尺寸比例图（自上而下、从左到右）：土星、木星、天王星、海王星（外行星），地球、金星、火星和水星（内行星）。} \label{fig_Planet_1}
\end{figure}
行星是一种巨大且呈圆形的天体，一般要求其绕恒星、恒星遗迹或褐矮星运行，并且其自身不属于这些天体。按该术语最严格的定义，太阳系共有八颗行星：类地行星（水星、金星、地球和火星）以及巨行星（木星、土星、天王星和海王星）。关于行星形成的最佳理论是星云假说，该假说认为一团星际云从星云中坍缩，形成一颗年轻的原恒星，而这颗原恒星被一个原行星盘所环绕。行星在该盘内通过由引力驱动的物质逐渐累积而成长，这一过程称为吸积。

“行星”（planet）一词源于希腊语 πλανήται（planḗtai），意为“流浪者”。在古代，这个词指的是太阳、月亮，以及在恒星背景上肉眼可见并会移动的五个光点——水星、金星、火星、木星和土星。行星在历史上具有宗教关联：多种文化将天体与神祇对应，而这种与神话和民间传说的联系，也延续到了对新发现的太阳系天体的命名方式之中。随着16至17世纪日心说取代地心说，地球本身也被认作一颗行星。

随着望远镜的发展，“行星”一词的含义被扩展，纳入了只能借助观测工具才能看见的天体：地球之外行星的卫星；冰巨行星天王星与海王星；谷神星及后来被确认属于小行星带的其他天体；以及冥王星，后来被发现是柯伊伯带中冰质天体族群中最大的成员。对柯伊伯带中其他大型天体的发现，尤其是厄里斯，引发了关于如何精确定义“行星”的争论。2006年，国际天文学联合会（IAU）通过了太阳系中“行星”的定义，将四颗类地行星与四颗巨行星归入行星范畴；而谷神星、冥王星与厄里斯则被归入“矮行星”类别。尽管如此，许多行星科学家仍持续以更宽泛的方式使用“行星”一词，包括矮行星以及如月球般呈圆形的卫星。

天文学的进一步发展使得人们发现了超过5900颗太阳系外的行星，即系外行星。这些系外行星常呈现太阳系行星所不具备的特殊性质，例如“热木星”——如飞马座51b般绕其母恒星极近距离运行的巨行星——以及极端偏心的轨道，如HD~20782~b。褐矮星与比木星更大的行星的发现也推动了对于定义边界的讨论，即如何划定行星与恒星之间的分界线。已有多颗系外行星被发现运行在其恒星的宜居带内（即行星表面可能存在液态水的区域），但地球仍是唯一已知能够承载生命的行星。宜居带的形成受多种因素影响，例如行星的大气组成（若其具有大气）。例如火星在赤道区域的白天温度足以形成液态水，只是其气压过低而无法实现。
\subsection{形成}
\begin{figure}[ht]
\centering
\includegraphics[width=8cm]{./figures/edc16278f48809b3.png}
\caption{原行星盘} \label{fig_Planet_2}
\end{figure}
\begin{figure}[ht]
\centering
\includegraphics[width=8cm]{./figures/645312c0cec1ae7e.png}
\caption{行星形成过程中原行星之间的碰撞} \label{fig_Planet_4}
\end{figure}
行星的形成机制目前尚未被完全确定。主流理论认为，行星在星云塌缩形成一层薄的气体和尘埃盘的过程中逐渐聚合而成。在中心区域形成原恒星，其周围则是一个旋转的原行星盘。通过吸积（即通过“粘性碰撞”的过程），盘中的尘埃粒子逐渐聚集质量，形成越来越大的天体。在局部区域中会出现称为微行星的质量集中体，它们通过自身的引力吸引更多物质，从而加速吸积过程。这些质量集中体变得越来越致密，最终在自引力作用下向内塌缩，形成原行星 \(^\text{[6]}\)。

当一个原行星的质量增长到略高于火星的质量时，它会开始吸积一层延展的大气层\(^\text{[7]}\)，从而通过大气阻力效应\(^\text{[8][9]}\)极大提高其对微行星的捕获效率。根据固体与气体的吸积历史不同，最终可能形成巨行星、冰巨行星或类地行星\(^\text{[10][11][12]}\)。人们认为木星、土星和天王星的规则卫星可能以类似方式形成\(^\text{[13][14]}\)；然而，海王星的卫星海卫一很可能是被捕获的\(^\text{[15]}\)。地球的月球\(^\text{[16]}\)与冥王星的卫星卡戎\(^\text{[17]}\)则可能是在巨型碰撞中形成。

当原恒星增大至能够点燃并成为真正的恒星时，残余的吸积盘会从内向外被光致蒸发、太阳风、Poynting–Robertson 阻力及其他效应逐渐清除\(^\text{[18][19]}\)。此后，系统中可能仍存在许多原行星，它们相互环绕或环绕恒星运行，但随着时间推移，其中许多会发生碰撞，或合并成更大的原行星，或为其他原行星提供可吸积的碎片\(^\text{[20]}\)。那些质量足够大的天体会吸尽其轨道邻域的大部分物质而成为行星。避免了碰撞的原行星可能通过引力捕获成为行星的天然卫星，或者留存在某些轨道带中，最终成为矮行星或小天体\(^\text{[21][22]}\)。
\begin{figure}[ht]
\centering
\includegraphics[width=8cm]{./figures/624858ff17708327.png}
\caption{超新星遗迹喷 ejecta 产生行星形成物质} \label{fig_Planet_5}
\end{figure}
较小的微行星的高能撞击（以及放射性衰变）会加热正在成长的行星，使其至少部分熔化。行星内部开始按照密度分异，密度更高的物质下沉至核心。\(^\text{[23]}\)较小的类地行星在这种聚积过程中会失去其大部分大气，但失去的气体可通过地幔逸气以及随后彗星的撞击所补充\(^\text{[24]}\)（更小的行星会通过各种逃逸机制而失去它们获得的任何大气\(^\text{[25]}\)）。

随着对太阳以外恒星的行星系统的发现与观测，人们得以对上述理论加以扩展、修正甚至替代。金属丰度——一个天文学术语，用来描述原子序数大于 2（即氦）的化学元素的丰度——似乎决定着恒星拥有行星的可能性。\(^\text{[26][27]}\)因此，与金属贫乏的第二星族（Population II）恒星相比，金属丰富的第一星族（Population I）恒星更可能拥有规模可观的行星系统。\(^\text{[28]}\)
\subsection{太阳系中的行星}
根据 IAU 的定义，太阳系中共有八颗行星（按距太阳距离递增排列）\(^\text{[2]}\)：Mercury、Venus、Earth、Mars、Jupiter、Saturn、Uranus 与 Neptune。Jupiter 质量最大，约为地球质量的 318 倍；而 Mercury 最小，仅为地球质量的 0.055 倍。\(^\text{[29]}\)

太阳系行星可依据组成成分划分为不同类别。类地行星（terrestrials）与 Earth 类似，其主体由岩石与金属构成：Mercury、Venus、Earth 与 Mars。Earth 是类地行星中体积最大的一个。\(^\text{[30]}\)

巨行星（giant planets）的质量显著大于类地行星：Jupiter、Saturn、Uranus 与 Neptune。\(^\text{[30]}\)它们在成分上与类地行星存在根本差异。气体巨行星（gas giants）Jupiter 与 Saturn 主要由氢与氦组成，是太阳系中最为巨大的行星。Saturn 的质量约为 Jupiter 的三分之一，即约 95 个地球质量。\(^\text{[31]}\)

冰巨行星（ice giants）Uranus 与 Neptune 则主要由低沸点物质组成，如 water、methane 与 ammonia，并拥有由 hydrogen 与 helium 构成的厚大气层。它们的质量显著低于气体巨行星，仅分别为 14 与 17 个地球质量。\(^\text{[31]}\)
\begin{figure}[ht]
\centering
\includegraphics[width=8cm]{./figures/f984108e351cc1ba.png}
\caption{太阳、行星、矮行星与卫星的尺寸按比例标注示意。各天体之间的距离不按比例绘制。小行星带位于火星与木星轨道之间，柯伊伯带位于海王星轨道之外。} \label{fig_Planet_6}
\end{figure}
\noindent 矮行星是达到引力塑形（即具有趋于球形的外形），但尚未清除其轨道附近其他天体的天体。按照距太阳平均距离递增排列，被天文学界普遍认可的矮行星包括：Ceres，Orcus，Pluto，Haumea，Quaoar，Makemake，Gonggong，Eris 与 Sedna。\(^\text{[32][33]}\)Ceres 位于 Mars 与 Jupiter 轨道之间的小行星带，是该区域最大的天体。其余八个均位于 Neptune 轨道之外。Orcus、Pluto、Haumea、Quaoar 与 Makemake 运行于 Kuiper 带——这是位于 Neptune 轨道之外的第二个太阳系小天体带。Gonggong 与 Eris 运行于散射盘（scattered disc），其位置更为外侧，且该区域不同于 Kuiper 带——会因 Neptune 的引力扰动而呈现轨道不稳定性。Sedna 是目前已知最大的游离天体（detached object），其轨道不会接近太阳到足以受任何经典行星的显著影响；其轨道成因仍在讨论之中。以上九个天体均与类地行星相似，拥有固体表面，但其成分主要是冰与岩石，而非岩石与金属。此外，它们全部小于 Mercury，其中 Pluto 是目前已知体积最大的矮行星，而 Eris 的质量最大。\(^\text{[34][35]}\)

至少有十九颗具行星质量的卫星（planetary-mass moons 或 satellite planets），即质量足以呈现椭球形状的卫星\(^\text{[4]}\)：

\begin{itemize}
    \item Earth 的一颗卫星：Moon；
    \item Jupiter 的四颗卫星：Io，Europa，Ganymede 与 Callisto；
    \item Saturn 的七颗卫星：Mimas，Enceladus，Tethys，Dione，Rhea，Titan 与 Iapetus；
    \item Uranus 的五颗卫星：Miranda，Ariel，Umbriel，Titania 与 Oberon；
    \item Neptune 的一颗卫星：Triton；
    \item Pluto 的一颗卫星：Charon。
\end{itemize}
Moon、Io 与 Europa 的成分与类地行星相似；其它卫星的成分则与矮行星类似，以冰与岩石为主，其中 Tethys 几乎完全由冰组成。由于其表面的冰层使内部难以观测，Europa 常被视为一颗冰质行星。\(^\text{[4][36]}\)

Ganymede 与 Titan 在半径上均大于 Mercury，而 Callisto 也几乎与之相当，但这三者的质量均远小于 Mercury。Mimas 是公认的几何意义上的行星（即达到地质行星标准）中最小的一个，其质量约为地球质量的六百万分之一；然而，还有许多比它更大的天体却可能并非地质行星（例如 Salacia）。\(^\text{[32]}\)
\subsection{Exoplanets}
\begin{figure}[ht]
\centering
\includegraphics[width=8cm]{./figures/85a4bfec6f369e50.png}
\caption{截至 2023 年 8 月的每年系外行星探测数量（数据来自 NASA 系外行星档案）\(^\text{[37]}\)} \label{fig_Planet_8}
\end{figure}
系外行星（exoplanet）指位于太阳系之外的行星。截至 2025 年 10 月 30 日，已经确认的系外行星共有 6{,}128 颗，分布在 4{,}584 个行星系统中，其中 1{,}017 个系统拥有不止一颗行星。\(^\text{[38]}\)已知的系外行星尺寸跨度极大，从比 Jupiter 大两倍的气态巨行星到仅略大于 Moon 的微小行星均有分布。对引力微透镜（gravitational microlensing）数据的分析表明，在 Milky Way 中，每颗恒星平均至少拥有 1.6 颗束缚行星。\(^\text{[39]}\)

1992 年初，无线电天文学家 Aleksander Wolszczan 与 Dale Frail 宣布发现两颗绕脉冲星 PSR 1257+12 运行的行星。\(^\text{[40]}\)这项发现随后得到证实，并普遍被认为是首次确定无疑的系外行星探测。研究者推测，这些行星可能形成于产生该脉冲星的超新星爆炸后残留的盘状物质中。\(^\text{[41]}\)

首次确认的、绕普通主序星运行的系外行星于 1995 年 10 月 6 日被发现。当时 Geneva 大学的 Michel Mayor 与 Didier Queloz 宣布在恒星 51 Pegasi 周围探测到行星 51 Pegasi b。\(^\text{[42]}\)从那时直到 Kepler 空间望远镜任务开始前，人类已知的大多数系外行星都是质量与 Jupiter 相当或更大的气态巨行星，因为它们更易被探测到。Kepler 候选行星目录中大多数行星的尺度介于 Neptune 与更小的行星之间，有些甚至小于 Mercury。\(^\text{[43][44]}\)

2011 年，Kepler 团队报告了首次发现绕类太阳恒星运行的地球大小行星：Kepler-20e 与 Kepler-20f。\(^\text{[45][46][47]}\)自那时起，已有超过 100 颗大小近似 Earth 的行星被确认，其中有 20 颗位于其恒星的宜居带（habitable zone）内——这是允许类地行星在具备足够大气压的条件下，其表面可能存在液态水的轨道范围。\(^\text{[48][49][50]}\)据估计，每五颗类太阳恒星中就有一颗拥有位于宜居带中的地球大小行星，这意味着距离地球最近的此类行星预计不会超过 12 光年。\(^\text{[a]}\)此类类地行星的出现频率是 Drake 方程中的变量之一，该方程用于估算 Milky Way 中存在的具备智慧并能够通信的文明数量。\(^\text{[53]}\)

在太阳系中不存在的行星类型包括超地球（super-Earth）与迷你海王星（mini-Neptune），其质量介于 Earth 与 Neptune 之间。质量不超过约两倍地球质量的天体通常被认为具有类似 Earth 的岩石组成；超过该范围，则会演化为含有大量挥发物与气体、类似 Neptune 的结构。\(^\text{[54]}\)行星 Gliese~581c 的质量为 Earth 的 $5.5$--$10.4$ 倍\(^\text{[55]}\)，在被发现时因可能处于其恒星的宜居带而引起关注\(^\text{[56]}\)，但后续研究表明它实际上距离其恒星过近，无法宜居。\(^\text{[57]}\)更大质量的行星亦已被发现，其性质与质量向上连续延伸至褐矮星（brown dwarf）领域。\(^\text{[58]}\)

已知一些系外行星与其母恒星之间的距离远小于太阳系中任意一颗行星与 Sun 的距离。太阳系内距离 Sun 最近的 Mercury 距离为 $0.4\ \mathrm{AU}$，其公转周期为 $88$ 天；但超短周期行星（ultra-short-period planets）可以在不到一天内完成公转。Kepler-11 系统中有五颗行星的轨道周期都短于 Mercury，而且它们的质量都远大于 Mercury。  
另有热木星（hot Jupiters），例如 51~Pegasi~b\(^\text{[42]}\)，紧贴恒星运行，可能因蒸发而变成“冥狱行星”（chthonian planets），也就是仅剩下残核的行星。  
此外，还有系外行星比太阳系中任何行星离其恒星更远。Neptune 距 Sun 为 $30$ AU，公转周期为 $165$ 年；但某些系外行星距离其恒星数千 AU，公转周期超过百万年（如 COCONUTS-2b）。\(^\text{[59]}\)
\subsection{属性}
尽管每颗行星都具有各自独特的物理特征，但它们之间仍然存在一些普遍的共同点。其中一些特征——如行星环或天然卫星——目前仅在太阳系行星中被观测到，而另一些则在系外行星中也普遍存在。\(^\text{[60]}\)
\subsection{动力学特征}
\textbf{轨道（Orbit）}
\begin{figure}[ht]
\centering
\includegraphics[width=8cm]{./figures/0d173cd711ffbba6.png}
\caption{行星海王星的轨道与冥王星轨道的比较。请注意冥王星轨道相对于海王星轨道的拉长（偏心率），以及其相对于黄道面的巨大夹角（轨道倾角）。} \label{fig_Planet_7}
\end{figure}
在太阳系中，所有行星绕日公转的方向都与太阳自转方向一致：从太阳北极上方看为逆时针方向。至少已经发现一颗系外行星（WASP-17b）其公转方向与其恒星自转方向相反。\(^\text{[61]}\)一个行星绕恒星运行一周的周期称为其恒星周期或“年”。\(^\text{[62]}\)行星的“年长”取决于它与恒星的距离：距离越远，轨道越长，同时由于恒星引力作用较弱，其运行速度也越慢。

没有任何行星的轨道是完全圆形的，因此其与恒星的距离会在一年内发生变化。行星最接近恒星的位置称为“近日点”（在太阳系中称“近日点”）；最远的位置称为“远日点”。当行星接近近日点时，其速度会增大，因为它将引力势能转化为动能，就像地球上自由落体物体下落时会加速一样。当地行星接近远日点时，其速度会减小，就像向上抛出的物体在接近最高点时速度会逐渐减小一样。\(^\text{[63]}\)

每个行星的轨道由一组轨道要素来描述：
\begin{itemize}
    \item 偏心率（eccentricity）描述行星椭圆（椭圆形）轨道的扁长程度。偏心率低的行星轨道更接近圆形，而偏心率高的行星轨道则更为椭圆。太阳系中的行星及大型卫星的偏心率都相对较低，因此它们的轨道几乎是圆形的。\(^\text{[62]}\)彗星、许多柯伊伯带天体以及若干系外行星则具有非常高的偏心率，因此其轨道极为椭圆。\(^\text{[64][65]}\)
    \item 半长轴（semi-major axis）给出轨道的尺度。它是椭圆轨道最长直径的一半。此值并不等同于远日点距离，因为没有任何行星的轨道是以恒星为椭圆几何中心的。\(^\text{[62]}\)
    \item 轨道倾角（inclination）表示行星轨道相对于某一参考平面的倾斜程度。在太阳系中，参考平面是地球轨道平面，称为黄道面。对于系外行星，参考平面称为“天空平面”，即垂直于地球观察方向的平面。\(^\text{[66]}\)太阳系八大行星的轨道都非常接近黄道面；然而，一些较小的天体，如智神星（Pallas）、冥王星（Pluto）和阋神星（Eris），以及彗星的轨道倾角都要大得多。\(^\text{[67]}\)大型卫星的轨道通常也不太偏离其母行星的赤道，但地球的月球、土星的土卫八（Iapetus）以及海王星的海卫一（Triton）是例外。其中海卫一尤为独特，它是唯一逆行绕行母行星的大型卫星，即其轨道方向与母行星自转方向相反。\(^\text{[68]}\)
    \item 行星轨道穿过参考平面的两个点称为升交点（ascending node）和降交点（descending node）。\(^\text{[62]}\)升交点经度（longitude of the ascending node）是从参考平面零经度到升交点的角度。近日点参数（argument of periapsis，太阳系中称近日点参数）是从升交点到行星最接近恒星位置之间的角度。\(^\text{[62]}\)
\end{itemize}
\textbf{自转轴倾角}
\begin{figure}[ht]
\centering
\includegraphics[width=8cm]{./figures/9e2bb5a76a5b4400.png}
\caption{地球的自转轴倾角约为 $23.4^\circ$。它在 $22.1^\circ$ 与 $24.5^\circ$ 之间以约 $41{,}000$ 年为周期振荡，目前正处于减小阶段。} \label{fig_Planet_9}
\end{figure}
行星具有不同程度的自转轴倾角；它们相对于其恒星赤道平面呈一定角度自转。这导致行星表面各半球在其公转周期内所接收的光照量发生变化：当北半球背向恒星时，南半球则朝向恒星，反之亦然。因此，每颗行星都会经历季节，其气候会随着一年中的位置变化而改变。每个半球在其自转轴最远离或最近指向恒星的位置，称为至点。每颗行星在一个公转周期中都有两个至点；当一个半球处于夏至且白昼时间最长时，另一半球则处于冬至且白昼最短。各半球所接收的光和热量的差异，会在每年引起对应半球天气模式的变化。

木星的自转轴倾角非常小，因此其季节变化极其微弱；而天王星的倾角则极端到几乎“侧躺”着自转，这意味着其半球在至点附近要么持续处于日照，要么持续处于黑夜。\(^\text{[69]}\)在太阳系中，水星、金星、谷神星和木星的倾角非常小；智神星、天王星和冥王星的倾角极大；而地球、火星、灶神星、土星和海王星的倾角属于中等范围。\(^\text{[70][71][72][73]}\)

对于系外行星，其自转轴倾角目前尚不确定，但由于靠近母恒星，大多数“热木星”被认为倾角极小。\(^\text{[74]}\)类似地，行星质量的卫星自转轴倾角通常接近零，\(^\text{[75]}\)其中地球的月球以 $6.687^\circ$ 为最大例外；\(^\text{[76]}\)此外，木卫四（卡利斯托）的自转轴倾角在数千年的时间尺度上会在 $0^\circ$ 到约 $2^\circ$ 之间变化。\(^\text{[77]}\)

\textbf{自转}

行星围绕穿过其中心的不可见自转轴进行自转。行星自转一周的时间称为恒星日（stellar~day）。太阳系中大多数行星的自转方向与它们绕太阳公转的方向一致，即从太阳北极上方观察为逆时针方向。然而，金星\(^\text{[78]}\)与天王星\(^\text{[79]}\)是例外，它们呈顺时针方向自转。不过，由于天王星自转轴倾角极端，对于其“北极”方向存在不同的约定，因此关于它是顺时针还是逆时针自转也存在不同定义。\(^\text{[80]}\)无论采用哪种约定，天王星相对于其轨道而言都属于逆行自转（retrograde~rotation）。\(^\text{[79]}\)
\begin{figure}[ht]
\centering
\includegraphics[width=8cm]{./figures/0c97e1c28cdd96c6.png}
\caption{行星与月球的自转周期（加速 10\,000 倍；负值表示逆行自转）、扁率以及自转轴倾角的比较（SVG 动画）} \label{fig_Planet_10}
\end{figure}
行星的自转可在形成过程中由多种因素诱发。吸积物体各自所携带的角动量会汇总为行星的总体角动量；巨行星吸积气体的过程也会增加其自转角动量；最后，在行星形成后期，原行星之间的随机吸积碰撞会随机改变行星的自转轴方向。\(^\text{[81]}\)

行星昼长差异极大：金星完成一次自转需要 243 天，而巨行星仅需数小时。\(^\text{[82]}\)系外行星的自转周期通常未知，但对“热木星”而言，由于其极度接近母恒星，它们会被潮汐锁定（即其自转与公转同步）。因此，它们始终以同一面朝向恒星，一侧永昼，一侧永夜。\(^\text{[83]}\) 水星与金星（最靠近日面的两颗行星）同样具有非常缓慢的自转：水星被潮汐锁定在 3:2 的自旋–轨道共振中（每绕太阳两周自转三次），\(^\text{[84]}\)而金星的自转可能处于潮汐力减速与由太阳加热驱动的大气潮汐加速之间的一种平衡状态。\(^\text{[85][86]}\)

所有的大型卫星都被潮汐锁定于其母行星；\(^\text{[87]}\)冥王星与其卫星卡戎互相潮汐锁定，\(^\text{[88]}\)妖神星（Eris）与其卫星迪斯诺米亚（Dysnomia）亦然，\(^\text{[89]}\)甚至可能包括阎王星（Orcus）与其卫星 Vanth。\(^\text{[90]}\)其他已知自转周期的矮行星自转速度都比地球快；妊神星（Haumea）的自转极快，以至被拉伸成一个三轴椭球体。\(^\text{[91]}\)系外行星 Tau Boötis b 与其母星 Tau Boötis 似乎也互相潮汐锁定。\(^\text{[92][93]}\)

\textbf{轨道清空}

根据国际天文学联合会（IAU）的定义，行星的关键动力学特征在于其已经清空了自身轨道邻域。一颗已清空邻域的行星，其质量足够大，能够聚集或清扫掉位于其轨道上的所有微行星（planetesimals）。实际上，它是独立围绕恒星运行的，而不是与大量尺寸相近的天体共享同一轨道。如上所述，这一特征被纳入 IAU 于 2006 年 8 月通过的行星正式定义中。[2] 虽然这一标准目前仅适用于太阳系，但已经发现若干年轻的系外行星系统，其原恒星盘（circumstellar disc）中存在能量或结构迹象，暗示轨道清空过程正在发生。[94]
\subsubsection{物理特征}
\textbf{大小与形状}

行星在自身重力作用下被拉向近似球形，因此行星的大小可以用平均半径来粗略表达（例如地球半径、木星半径）。然而，行星并非完美球体；例如，由于地球自转，其两极略微扁平，赤道附近略有鼓起。[95] 因此，用扁球体（oblate spheroid）来近似地球形状更为准确，其赤道直径比极轴直径大约 43 公里（27 英里）。[96]

一般而言，行星的形状可以通过给出扁球体的极半径与赤道半径，或指定一个参考椭球体来描述。由此可以计算行星的扁率、表面积和体积；进一步结合其大小、形状、自转速率以及质量，即可计算其正常重力。[97]

\textbf{质量}

行星最重要的物理特征在于：其质量足够大，使自身重力能够压倒维系其物理结构的电磁力，使其达到\textbf{流体静力学平衡（hydrostatic equilibrium）}。这意味着所有行星都是球形或球状体。低于某一质量上限时，天体可以保持不规则形状；但当其化学组成所允许的临界质量被超过时，重力会开始将物体拉向自身质心，使其最终塌缩成球体。[98]

质量也是区分行星与恒星的主要属性。太阳与木星之间的质量区间在太阳系中无任何天体，但系外行星中存在此类天体。恒星的最低质量极限约为 75–80 个木星质量（MJ）。部分学者提议把这一界限作为行星的最大质量上限，因为从“显著自压缩开始”（约一土星质量）一直到氢点火成为红矮星，其内部物理并无本质变化。[54]另一些观点认为约 13 MJ（太阳丰度条件下）可以作为行星与褐矮星的分界，因为在此质量上天体能够发生氘核聚变。[99] 然而氘燃烧持续时间短，大多数褐矮星早已结束其氘燃烧阶段。[58] 这一界限并未得到普遍认可：Exoplanets Encyclopaedia 将上限取到 60 MJ，Exoplanet Data Explorer 则取到 24 MJ。[100][101]

已知质量最小、质量测量可靠的系外行星是 1992 年发现的脉冲星行星 PSR B1257+12A，其质量约为水星的一半。[102] 更小的是围绕白矮星运行的 WD 1145+017 b，其质量约与矮行星妊神星（Haumea）相当，通常被称为小行星（minor planet）。[103]围绕主序星（非太阳）运转、质量最小的已知行星是 Kepler-37b，其质量（和半径）可能略大于月球。[44] 太阳系中普遍认可的最小地球物理行星是土星的卫星米玛斯（Mimas），半径约为地球的 3.1\%，质量约为地球的 0.00063\%。[104] 土星的另一颗更小的卫星——福柏（Phoebe），目前是不规则天体，半径为地球的 1.7\%，质量为 0.00014\%，但据信其在早期曾达到流体静力平衡并发生内部分异，之后被撞击破坏而变形。[105][104][106]部分小行星可能是开始聚积并发生分异的原行星（protoplanets）的碎块，但因灾难性碰撞，仅存金属核或岩质核心；[107][108][109] 或者是这些碰撞碎片再聚集形成的天体。[110]

\textbf{内部分异}
\begin{figure}[ht]
\centering
\includegraphics[width=8cm]{./figures/a7021b87c5306863.png}
\caption{木星内部结构示意图：在岩石核心之上覆盖着深厚的金属氢层} \label{fig_Planet_11}
\end{figure}
每一颗行星在起初形成时都处于完全流体状态；在形成早期，更致密、更重的物质下沉至中心，使较轻的物质留在靠近表面的区域。因此，每颗行星都拥有分异的内部结构，由致密的行星核心组成，外围包围着一个曾经或仍然处于流体状态的地幔。类地行星的地幔被坚硬的地壳封闭在内部[111]，但在巨行星中，地幔与其上方的云层逐渐过渡并融合。类地行星的核心主要由铁和镍等元素构成，其地幔则主要由硅酸盐组成。木星和土星被认为拥有由岩石与金属构成的核心，外包一层金属氢地幔[112]。天王星和海王星体积较小，它们拥有岩石核心，外围包围着由水、氨、甲烷及其他冰类物质构成的地幔[113]。这些行星核心内部的流体运动产生了地磁发电机效应，从而生成行星的磁场[111]。

类似的内部分异过程也被认为发生在部分大型卫星与矮行星上[32]，但这一过程并非总能完全进行：谷神星（Ceres）、卡利斯托（Callisto）以及泰坦（Titan）似乎都仅完成了部分分异[114][115]。小行星灶神星（Vesta）虽然因受到撞击破坏而失去球形、未被归类为矮行星，但其内部结构已经完成分异[116]，与金星、地球和火星的情况相似[109]。

\textbf{大气}
\begin{figure}[ht]
\centering
\includegraphics[width=8cm]{./figures/01c321ee2d5369eb.png}
\caption{“地球大气层”} \label{fig_Planet_12}
\end{figure}
太阳系中除水星以外的所有行星都拥有相当稠密的大气层[117]，因为它们的引力足够强，可以将气体束缚在靠近表面的区域。土星最大的卫星——泰坦（Titan）也拥有稠密的大气层，其厚度甚至超过地球大气层[118]；海王星最大的卫星——海卫一（Triton）[119]以及矮行星冥王星则拥有更为稀薄的大气层[120]。体积更大的巨行星质量足以维持大量轻质气体（如氢和氦），而体积较小的行星则会将这些气体流失至太空[121]。对系外行星的分析表明，保留这些轻质气体的临界质量约为 $2.0^{+0.7}_{-0.6}\,M_{\oplus}$，因此地球与金星大致位于岩质行星大小的上限附近[54]。

地球大气的成分与其他行星不同，因为生命过程向大气中引入了游离的分子氧[122]。火星与金星的大气都以二氧化碳为主，但其密度差异极大：火星平均表面气压不到地球的 1\%（不足以允许液态水存在）[123]，而金星平均表面气压约为地球的 92 倍[124]。金星目前极端的高温高压环境被认为源于其历史上的失控温室效应，使其成为太阳系中表面温度最高的行星，甚至超过更靠近太阳的水星[125]。尽管金星表面条件极端严酷，其大气中约 50–55 公里高度的温度与压力却与地球相近（这是太阳系中除地球外唯一存在此类条件的区域），因而该高度区域被提出为未来人类探索的可能基地[126]。泰坦拥有除地球外太阳系中唯一富含氮气的稠密行星大气。类似于地球表面条件接近水的三相点，使水能在三种物态间循环；泰坦的环境条件则接近甲烷的三相点[127]。

行星大气在不同程度上受入射辐射或内部能量驱动，导致其形成动态的天气系统，例如：地球上的飓风、火星上的全球性沙尘暴、木星上比地球还大的反气旋（大红斑），以及海王星大气中的空洞[69]。在系外行星上观察到的天气现象包括：在 HD 189733 b 上出现的、面积为木星大红斑两倍的高温区域[128]；以及位于热木星 Kepler-7b、超级地球 Gliese 1214 b 等行星上的云层[129][130][131]。

由于热木星与其宿主恒星极端接近，它们的大气在恒星辐射下不断被剥离，如同彗星般形成尾迹[132][133]。这类行星常在昼夜两侧呈现巨大温差，从而可能产生超音速风流[134]；然而，影响昼夜温差与大气动力学的因素众多，其细节十分复杂[135][136]。

\textbf{磁层}
\begin{figure}[ht]
\centering
\includegraphics[width=8cm]{./figures/19b83eb33522eada.png}
\caption{地球磁层（示意图）} \label{fig_Planet_13}
\end{figure}
行星的一项重要特征是其本征磁矩，而本征磁矩又会产生磁层。磁场的存在表明行星在地质学意义上仍然“活着”。换言之，具有磁化性质的行星在其内部存在电导流体的流动，从而产生磁场。这些磁场会显著改变行星与太阳风之间的相互作用。一个具有磁化的行星会在其周围的太阳风中开辟出一个空腔，这个区域称为磁层（magnetosphere），太阳风无法穿透其中。磁层的尺寸可能远大于行星本身。相反，非磁化行星只能依靠电离层与太阳风的相互作用而形成较小的诱导磁层，这种磁层无法有效保护行星[137]。

在太阳系八大行星中，只有金星和火星缺乏这样的磁场[137]。在具有磁场的行星中，水星的磁场最弱，几乎无法偏转太阳风。木星的卫星伽利略卫星之一——盖尼米德（Ganymede）拥有更强的磁场，而木星自身的磁场是太阳系最强的（其强度之大，事实上对未来所有前往其轨道内卡利斯托以内的载人任务构成严重的健康风险[138]）。其他巨行星在其表面的磁场强度与地球相近，但它们的磁矩却显著更大。天王星和海王星的磁场相对于其自转轴具有强烈的倾角，并且磁场中心相对于行星中心存在明显偏移[137]。

2003年，夏威夷的一组天文学家在观测恒星 HD 179949 时发现其表面出现一个亮斑，显然是由其轨道上一颗热木星的磁层所诱发的[139][140]。
\subsubsection{次级特征}
\begin{figure}[ht]
\centering
\includegraphics[width=8cm]{./figures/7eabc59da7d0513c.png}
\caption{土星的行星环} \label{fig_Planet_15}
\end{figure}
在太阳系中，若干行星或矮行星（例如海王星和冥王星）具有彼此之间，或与更小天体之间的轨道周期共振现象。这种情况在卫星系统中尤为常见（例如木星的木卫一、木卫二与木卫三之间的共振，或土星的土卫二与土卫四之间的共振）。除水星与金星之外，所有行星都拥有天然卫星（通常称为“月亮”）。地球有一颗卫星，火星有两颗，而巨行星拥有数量众多、结构复杂的“类行星系统”般的卫星群。除谷神星（Ceres）与赛德娜（Sedna）外，所有共识矮行星也都至少拥有一颗卫星。许多巨行星的卫星与类地行星及矮行星具有相似的地质特征，其中一些被研究为可能适宜生命存在的地点（尤其是木卫二与土卫二）[141][142][143][144][145]。

四大巨行星都拥有不同规模和复杂度的行星环。这些环主要由尘埃或颗粒组成，但可能包含微小的“卫星状微体”（moonlets），其引力有助于塑造并维持环的结构。虽然行星环的确切起源尚未明晰，但普遍认为它们源自掉入母行星洛希极限（Roche limit）内部、因潮汐力而被撕裂的天然卫星[146][147]。矮行星妊神星（Haumea）[148]和夸欧尔（Quaoar）[149]也被发现拥有行星环。

在系外行星周围尚未观测到类似的次级结构。亚褐矮星 Cha 110913−773444（常被描述为“流浪行星”）被认为拥有一个微小的原行星盘\,[150]，而亚褐矮星 OTS 44 被证实拥有质量至少相当于十个地球质量的巨大原行星盘\,[151]。
\subsection{历史与词源}
行星的概念在天文学历史上不断演变，从古代被视为“神圣之光”的天体，发展为科学时代的类地天体。该概念也已从仅指太阳系内的世界扩展至无数系外行星系统。关于“哪些天体应当被视为行星”这一共识曾多次改变：历史上，行星的范畴曾包括小行星、卫星以及像冥王星这样的矮行星[152][153][154]，而时至今日，这一问题仍未完全达成一致[154]。
\subsubsection{古代文明与古典行星}
由于太阳系中的五颗古典行星可被肉眼直接观测到，它们自古以来便为人类所知，并对神话、宗教宇宙观以及古代天文学产生了深远影响。古代天文学家注意到，某些光点会在天空中移动，与保持相对位置恒定的“恒星”不同。[155] 古希腊人将这些移动的光点称为 πλάνητες ἀστέρες（planētes asteres，“游移之星”），或简称为 πλανῆται（planētai，“游者”），[156] 今日的术语“planet（行星）”便源自此。[157][158][159]在古希腊、中国、巴比伦，以及几乎所有前现代文明中，[160][161] 人们普遍相信地球是宇宙的中心，所有的“行星”都绕地球运行。形成这一认知的原因包括：恒星与行星似乎每天都绕着地球旋转，[162] 以及基于常识的直观判断——地球坚实而稳定，没有运动迹象，似乎处于静止状态。[163]

\textbf{巴比伦}

已知最早拥有成熟行星理论的文明是巴比伦人，他们在公元前第一与第二千纪居住于美索不达米亚。现存最古老的行星天文文本为《阿米萨杜卡的金星泥板》（\emph{Babylonian Venus tablet of Ammisaduqa}），这是一份公元前7世纪的抄本，记录了对金星运动的观测，其原本可能早至公元前第二千纪。[164]《MUL.APIN》是两块源自公元前7世纪的楔形文字泥板，系统展示了太阳、月亮与行星在一年中的运行规律。[165]巴比伦后期的天文学是西方天文学乃至整个西方精确科学传统的起源。[166]《Enuma anu enlil》，写于新亚述时期（公元前7世纪），[167] 包含了一份凶兆列表以及它们与各种天象（包括行星运动）的对应关系。[168][169] 巴比伦天文学家识别了内行星金星、 水星，以及外行星火星、木星、土星，这些在望远镜发明前一直是人类已知的全部行星。[170]

\textbf{希腊—罗马天文学}

古希腊人最初并不像巴比伦人那样强调行星的重要性。公元前6至5世纪，毕达哥拉斯学派似乎发展了独立的行星体系，其中地球、太阳、月亮与行星都围绕宇宙中心的“中央之火”（\emph{Central Fire}）运行。据说毕达哥拉斯或巴门尼德斯最早意识到昏星（\emph{Hesperos}）与晨星（\emph{Phosphoros}）实际上是同一颗星（阿佛洛狄忒，即后来的金星），[171] 尽管这一点在美索不达米亚早已为人所知。[172][173] 公元前3世纪，萨摩斯的阿里斯塔克斯提出了日心体系，认为地球与行星绕太阳运行。然而地心体系一直占主导地位直至科学革命时期。[163]

到公元前1世纪的希腊化时期，希腊人开始发展自己的数学方法来预测行星位置。这些基于几何学而非巴比伦人算术方法的理论，最终在复杂度与完备性上超越了巴比伦体系，能够解释肉眼可见的绝大部分天体运动。该体系的最高成就在公元2世纪托勒密的《天文学大成》（\emph{Almagest}）中实现，其影响力巨大，以至于取代了所有更早的天文学著作，并在西方世界作为权威天文文本持续了13个世纪。[164][174]对希腊人与罗马人而言，可见的“行星”共有七个，它们都被认为依照托勒密定义的复杂运动规律绕地球运行。按托勒密的顺序并使用现代名称，这七个天体依次是：月球、水星、金星、太阳、火星、木星与土星。[159][174][175]
\subsubsection{中世纪天文学}
\begin{figure}[ht]
\centering
\includegraphics[width=8cm]{./figures/200a04b44cceda48.png}
\caption{1660 年克劳狄乌斯·托勒密地心模型的插图} \label{fig_Planet_16}
\end{figure}
在西罗马帝国灭亡之后，天文学在印度和中世纪伊斯兰世界中继续发展。公元 499 年，印度天文学家阿耶波多（Aryabhata）提出了一个行星模型，其中明确包含地球绕自身轴心的自转，并将这种自转解释为恒星看似向西移动的原因。他还推测行星运行轨道为椭圆形 [176]。阿耶波多的追随者在印度南部尤为兴盛，他们遵循他的地球自转等原理，并基于这些思想撰写了大量后续著作 [177]。

伊斯兰黄金时代的天文学主要在中东、中亚、安达卢斯与北非开展，之后传播至远东与印度。这些天文学家（如多才多艺的伊本·海赛姆 Ibn~al-Haytham）通常接受地心宇宙模型，但他们质疑托勒密的本轮体系并试图提出替代理论。10 世纪的天文学家阿布·赛义德·西吉（Abu~Sa'id~al-Sijzi）接受了地球绕自转轴自转的观点 [178]。11 世纪时，伊本·西那（Avicenna）观测到了金星凌日现象 [179]。其同时代学者比鲁尼（Al-Biruni）利用三角学提出了一种测量地球半径的新方法，与埃拉托色尼的方法不同，这一方法只需在一座高山进行观测即可 [180]。
\subsubsection{科学革命与外行星的发现}
\begin{figure}[ht]
\centering
\includegraphics[width=8cm]{./figures/43a174f81fae1e03.png}
\caption{一幅由 Emanuel Bowen 于 1747 年制作的真实比例太阳系海报。当时，天王星、海王星以及小行星带都尚未被发现。} \label{fig_Planet_18}
\end{figure}
随着科学革命的兴起以及哥白尼、伽利略和开普勒的日心模型的确立，“行星”一词的使用发生了改变：其含义不再是相对于恒星背景在天空中移动的天体，而是指直接（一级行星）或间接（二级行星或卫星行星）绕太阳运行的天体。由此，地球被加入到行星的名单之中 [181]，而太阳被移除。哥白尼体系中的一级行星数量一直保持不变，直到 1781 年威廉·赫歇尔发现天王星 [182]。

当 17 世纪发现木星的四颗卫星（伽利略卫星）以及土星的五颗卫星时，它们与地球的月球一起，被归为绕一级行星运行的“卫星行星”或“二级行星”；不过在随后的几十年里，它们逐渐被简称为“卫星”。直到 20 世纪 20 年代之前，科学家普遍仍将行星的卫星也视为行星，不过这种用法在非科学界并不常见 [154]。

在 19 世纪的第一个十年中，又有四个新的“行星”被发现：谷神星（1801 年）、智神星（1802 年）、婚神星（1804 年）和灶神星（1807 年）。但很快人们意识到，它们与先前已知的行星有明显差异：它们共享同一区域——火星与木星之间的小行星带——并具有部分重叠的轨道。这个区域原本只预期存在一颗行星，而且这些天体的体积远小于其他所有行星；甚至有人怀疑它们可能是某颗更大行星破碎后的碎片。赫歇尔称它们为“小行星”（意为“类星体”），因为即便在最大型的望远镜中，它们也呈现出类似恒星的外观，没有可分辨的盘状结构 [153][183]。

这一局面持续稳定了四十年，但在 1840 年代，又有数颗新的小行星被发现（1845 年的 Astraea；1847 年的 Hebe、Iris 和 Flora；1848 年的 Metis；以及 1849 年的 Hygiea）。新的“行星”几乎每年都有发现；因此，天文学家开始将这些小行星（又称小行星、次级行星）与主要行星分开编目，并为它们分配编号，而不是抽象的行星符号 [153]，尽管它们仍然被视为小型行星 [184]。

海王星于 1846 年被发现，其位置得以预测，是由于它对天王星的引力影响。然而，由于水星的轨道似乎也出现了类似的扰动，19 世纪末人们曾经认为太阳附近可能还存在另一颗行星。然而，水星轨道与牛顿引力预测之间的差异最终由爱因斯坦的广义相对论解释，而不是由另一颗行星解释 [185][186]。

冥王星于 1930 年被发现。由于最初的观测认为它比地球大 [187]，该天体立即被接受为第九颗主要行星。后来更精确的观测发现该天体实际上小得多：1936 年，Ray Lyttleton 提出冥王星可能是海王星的逃逸卫星 [188]；而 Fred Whipple 在 1964 年提出冥王星可能是一颗彗星 [189]。1978 年其大型卫星卡戎（Charon）的发现显示冥王星的质量仅为地球的 0.2\% [190]。由于这仍然远大于任何已知小行星，并且当时尚未发现其他海王星外天体，冥王星保留了行星地位，直到 2006 年才正式失去该地位 [191][192]。

20 世纪 50 年代，Gerard Kuiper 发表了关于小行星起源的论文。他认识到小行星通常并非如此前认为的那样呈球形，并指出小行星族群是碰撞的残余物。因此，他将最大的小行星视为“真正的行星”，而将较小的视为碰撞碎片。从 20 世纪 60 年代开始，“小行星”（asteroid）一词逐渐取代“次级行星”（minor planet），文献中将小行星视为行星的说法也变得罕见，除非讨论地质演化明显的三颗最大天体：谷神星、以及较少被提及的智神星和灶神星 [184]。

20 世纪 60 年代以来，随着空间探测器对太阳系的探索开始，行星科学重新受到重视。此时关于卫星定义的分歧出现：行星科学家开始重新考虑将大型卫星视为行星，而非行星科学的天文学家则通常不这样看待 [154]。（这与前一个世纪的定义并不完全一致，当时所有卫星都被视为次级行星，即使是如土星的 Hyperion 或火星的 Phobos、Deimos 这样非球形的天体也包括在内。）[193][194] 自那以来，八颗主要行星及其行星质量的卫星均已被航天器探测；许多小行星以及矮行星谷神星和冥王星也已被探测。然而，迄今为止，唯一被人类亲自探访的地球以外的行星质量天体仍然是月球。[b]
\subsubsection{界定“行星”一词}
越来越多的天文学家主张将冥王星取消行星地位，因为在 1990 年代和 21 世纪初，在太阳系同一区域（柯伊伯带）发现了许多与其大小相近的天体。事实证明，冥王星只是该区域数千个“小”天体之一 [195]。他们经常引用小行星降格为先例，尽管小行星被降格的原因主要基于其地质物理性质与真正行星的差异，而非因为它们位于小行星带中 [154]。一些较大的海王星外天体，如 Quaoar、Sedna、Eris 和 Haumea，[196] 曾被大众媒体誉为“第十颗行星”。

2005 年 Eris 的发现——其质量比冥王星大 27\%——成为制定行星官方定义的直接动力 [195]。如果将冥王星视为行星，那么逻辑上也必须将 Eris 视为行星。由于行星与非行星的命名程序完全不同，这造成了紧迫的局面：按照当时的规则，在没有界定“行星”的前提下，Eris 是不能被正式命名的 [154]。当时，人们还认为海王星外天体要达到“由自身引力变圆”的尺寸，与巨行星的卫星相近（直径约 400,km）。按照这一标准，柯伊伯带内可能有约 200 个呈圆形的天体，且在更远处可能有数千个 [197][198]。许多天文学家认为，公众不会接受一个导致行星数量剧烈膨胀的定义 [154]。

为解决这一问题，国际天文学联合会（IAU）着手制定“行星”的定义，并于 2006 年 8 月给出了正式版本。在这一定义下，太阳系被认为拥有八颗行星（即水星、金星、地球、火星、木星、土星、天王星和海王星）。那些满足前两条条件但不满足第三条的天体，被归类为矮行星，只要它们不是其他行星的天然卫星。最初，IAU 的一个委员会曾提出一个不包含条件 (c) 的定义方案，这一方案本会将更多的天体纳入行星范畴 [199]。但经过大量讨论，最终通过投票决定，这些天体应被归类为矮行星，而不是行星 [192][200]。

\textbf{对 IAU 定义的批评与替代方案}
\begin{figure}[ht]
\centering
\includegraphics[width=8cm]{./figures/88fe4ac34cbda9b1.png}
\caption{} \label{fig_Planet_17}
\end{figure}


